\documentclass[11pt,a4paper]{article}
\usepackage[margin=1in]{geometry}
\usepackage[T1]{fontenc}
\usepackage{lmodern,amsmath,amssymb,amsthm,booktabs}
\usepackage[hidelinks,pdftitle={A radix-seven counterexample to the cyclic-value conjecture},pdfauthor={Artem Dikarev, ChatGPT}]{hyperref}
\newtheorem{theorem}{Theorem}
\newtheorem{lemma}{Lemma}
\newcommand{\val}{\operatorname{val}}
\newcommand{\Z}{\mathbb Z}
\title{A radix-seven counterexample to Pratt's cyclic-value conjecture}
\author{Artem Dikarev \and ChatGPT}
\date{25 November 2026\\\small Preprint; independent review and priority verification pending}
\begin{document}
\maketitle
\begin{abstract}
Pratt conjectured that the number of zero-sum triples in equilateral trapezoid-free (ETF) triples of subsets of a cyclic group of order $N$ is at most $O(N^{1+\varepsilon})$ for every $\varepsilon>0$. We give a counterexample. For every positive integer $m$, with $k=8m$ and $N_k=3(7^k-1)+1$, we obtain $\val(\Z_{N_k})\ge (8m)!/(m!)^8$, giving a lower growth exponent of $\log_7 8>1$. The idea is to encode eight selected triples using base seven. Fixing the digit counts eliminates carries and, through a simple additive weight, excludes three unwanted local triples. A companion Lean proof establishes an alternative lower bound that also contradicts the conjecture.
\end{abstract}
\section{Introduction and statement}
How many zero-sum triples can three subsets of a cyclic group contain if certain overlapping configurations are forbidden? This is one of the questions raised by Pratt~\cite{pratt} in his study of generalized corners and matrix multiplication. We recall the precise restriction before describing our construction.

For three subsets $A,B,C$ of an additive finite abelian group $G$, write $R=\{(a,b,c)\in A\times B\times C:a+b+c=0\}$. The triple is equilateral trapezoid-free (ETF) if, for every fixed choice of anchors $a'\in A$, $b'\in B$ and $c'\in C$, each system
\[
\begin{aligned}
0&=a'+b+c=a+b'+c,\\
0&=a'+b+c=a+b+c',\\
0&=a+b'+c=a+b+c'
\end{aligned}
\]
has at most one solution $(a,b,c)\in A\times B\times C$. The quantity $\val(G)$ is the maximum of $|R|$ over such triples (Definition~17). Conjecture~24 asserts that $\val(\Z_N)=O(N^{1+\varepsilon})$ for every $\varepsilon>0$. Thus the conjecture rules out any fixed power improvement over linear growth.
\begin{theorem}\label{thm:main}
For every positive integer $m$, let $k=8m$ and $N_k=3(7^k-1)+1$. Then
\[
\val(\Z_{N_k})\ge \frac{(8m)!}{(m!)^8}.
\]
Consequently, along this infinite sequence of moduli,
\[
\val(\Z_{N_k})=\Omega\left(\frac{N_k^{\log_7 8}}{(\log N_k)^{7/2}}\right),
\qquad \log_7 8=1.06862156\ldots .
\]
\end{theorem}
The reason for using base seven is simple: eight choices per digit position grow faster than the range of numbers available in base-seven notation. If we could make those choices independently while preserving ETF, we would obtain roughly $8^k$ solutions on a range of size $7^k$. There is a difficulty, however. The sets of encoded numbers determine all their zero-sum triples; we cannot simply declare some of these triples forbidden. Our construction overcomes this by restricting the digit counts of each number. This costs only a polynomial factor in the number of solutions.

\section{An eight-edge digit relation}
We start with digit alphabets $A=\{0,1,2,4\}$, $B=\{0,1,\ldots,6\}$ and $C=\{0,1,5,6\}$. Among the triples in $A\times B\times C$ summing to six, let $Y$ denote the full relation and choose the following eight triples as our desired local relation:
\[
\begin{aligned}
X=\{&(0,0,6),(0,5,1),(0,6,0),(1,0,5),\\
     &(2,3,1),(2,4,0),(4,1,1),(4,2,0)\}.
\end{aligned}
\]
The other three elements of $Y$ are $(0,1,5),(1,4,1),(1,5,0)$. The full relation $Y$ is not ETF: for example, the anchors $a'=0$, $b'=1$ give two solutions $(0,1,5)$ and $(4,5,1)$ to the first anchored system with target sum six. We will retain the eight triples in $X$ and remove the three extras when we pass to longer words.
\begin{lemma}\label{lem:seed}
The abstract relation $X$ satisfies all three anchored uniqueness conditions. There exist integer functions $u,v,w$ on the digit alphabets such that $u(a)+v(b)+w(c)$ is zero on $X$ and one on $Y\setminus X$.
\end{lemma}
\begin{proof}
Since every edge sums to six, any two coordinates determine the third. Consider the first anchored system, with $a'$ and $b'$ fixed. A possible solution with third coordinate $c$ exists precisely when $a'$ and $b'$ both occur among the respective neighbors of $c$. Once $c$ is known, the sum equation determines the two remaining unknowns. Thus it suffices to show that the Cartesian products of the two neighbor sets of distinct $c$ are disjoint. The same argument applies with either of the other roles as the center. The following tables give the neighbor sets needed for all three checks.
\begin{center}
\begin{tabular}{c ll@{\qquad}c ll@{\qquad}c ll}
\toprule
$a$ & $B$ & $C$ & $b$ & $A$ & $C$ & $c$ & $A$ & $B$\\
\midrule
0 & 0,5,6 & 0,1,6 & 0 & 0,1 & 5,6 & 0 & 0,2,4 & 2,4,6\\
1 & 0 & 5 & 1 & 4 & 1 & 1 & 0,2,4 & 1,3,5\\
2 & 3,4 & 0,1 & 2 & 4 & 0 & 5 & 1 & 0\\
4 & 1,2 & 0,1 & 3 & 2 & 1 & 6 & 0 & 0\\
  & & & 4 & 2 & 0 & & &\\
  & & & 5 & 0 & 1 & & &\\
  & & & 6 & 0 & 0 & & &\\
\bottomrule
\end{tabular}
\end{center}
Within each of the three tables, the rectangles for distinct centers are disjoint. This proves the claimed uniqueness.

To remove the three extra triples later, we need a weight that distinguishes them using only the individual digits. Set $u(1)=1$, $v(0)=-2$, $w(5)=1$, $w(6)=2$, with all remaining values zero. The two selected triples with $b=0$ balance the negative weight $-2$ exactly, and the other six selected triples have weight zero. Each of the three extras has weight one. Thus $u(a)+v(b)+w(c)$ has the required values on all eleven triples.
\end{proof}
The key feature of this weight is its separability: its sum over a sequence of triples can be computed from the three digit words individually, without knowing how their digits are paired.
\section{Positional encoding and the induced relation}
Take $m\ge1$, $k=8m$ and $h=7^k-1$. Imagine writing each of the eight triples in $X$ exactly $m$ times, in any order, and reading the three columns as base-seven numbers. Rather than keeping the entire list of triples in the definition of our sets, we retain only the digit counts in each column. This gives three independently defined sets $A_k,B_k,C_k$ of encoded words of length $k$, with the following multiplicities (each count is multiplied by $m$):
\begin{center}
\begin{tabular}{ll}
\toprule
Role & Digit multiplicities\\
\midrule
$A$ & $0:3,\ 1:1,\ 2:2,\ 4:2$\\
$B$ & $0:2,\ 1:1,\ 2:1,\ 3:1,\ 4:1,\ 5:1,\ 6:1$\\
$C$ & $0:3,\ 1:3,\ 5:1,\ 6:1$\\
\bottomrule
\end{tabular}
\end{center}
The encoding $(d_i)_{i=0}^{k-1}\mapsto\sum_i d_i7^i$ is injective at fixed length and its image lies in $[0,h]$. The intended target sum is $h$, whose base-seven digits are all six. Our next task is to show that these separately defined sets cannot produce any unintended solutions.
\begin{lemma}\label{lem:filter}
Every induced integer target $a+b+c=h$ in these three sets consists coordinatewise of edges of $X$.
\end{lemma}
\begin{proof}
There are two possible sources of unwanted solutions: carries between positions, and extra triples within a position. We rule them out in that order. Suppose $a+b+c=h$. The prescribed counts fix the digit sums of the three words at $13m,21m,14m$, totaling $6k$. Ordinary base-seven addition gives nonnegative carries $q_0=q_k=0$ satisfying
\[
a_i+b_i+c_i+q_i=6+7q_{i+1}.
\]
Summing over $i$ gives $6\sum_{i=1}^{k-1}q_i=0$. In other words, the fixed digit sums leave no room for even a single nonzero carry. Every coordinate triple therefore belongs to $Y$.

It remains to exclude the three extra triples of $Y$. The prescribed counts force the sum of the weights from Lemma~\ref{lem:seed} over all positions to be $m-4m+3m=0$. Each position now contributes either zero (an edge of $X$) or one (an extra edge). Since no negative contribution is possible, the total can be zero only if every position belongs to $X$.
\end{proof}
\begin{lemma}\label{lem:count}
The integer relation at height $h$ is ETF and has exactly $(8m)!/(m!)^8$ targets.
\end{lemma}
\begin{proof}
Fix the anchors in any of the three systems, with target sum $h$. Both equations are target equations between elements of our sets, so Lemma~\ref{lem:filter} applies to each of them. At each position, the anchor digits are fixed and Lemma~\ref{lem:seed} allows at most one solution in $X$. The full words are therefore uniquely determined whenever a solution exists. This proves ETF in all three orientations.

We can now count solutions without worrying about additional ones. Every ordering of $m$ copies of each of the eight selected triples gives a target with the required digit counts. Conversely, by Lemma~\ref{lem:filter}, every target comes from a word over $X$. Each of the six digits $1,\ldots,6$ in the second role occurs in exactly one selected edge, so their counts force six edge multiplicities to equal $m$. The digits 5 and 6 in the third role fix the remaining two in the same way. The admissible edge words are exactly the permutations of this multiset. The encoding is injective, so their multinomial count is also the number of distinct targets.
\end{proof}
\section{Cyclic embedding and asymptotics}
\begin{proof}[Proof of Theorem~\ref{thm:main}]
To turn the integer construction into a cyclic one, we need to preserve the target equations without introducing modular wraparound. We take $N_k=3h+1$, reduce the first two sets modulo $N_k$, and translate the third by $-h$ before reducing. All three maps are injective since the original sets lie in $[0,h]$. For $a,b,c\in[0,h]$, the quantity $a+b+c-h$ belongs to $[-h,2h]$, where zero is the only multiple of $N_k$. Hence modular target equations are equivalent to their integer counterparts, in particular in all three anchored systems. Lemma~\ref{lem:count} gives the claimed ETF witness. Stirling's formula yields $(8m)!/(m!)^8=\Theta(8^k/k^{7/2})$, whereas $N_k=\Theta(7^k)$. The eight choices per position have thus survived at exponential scale, with only a polynomial loss, while the modulus grows like $7^k$. For any fixed $0<\varepsilon<\log_7 8-1$, the ratio $\val(\Z_{N_k})/N_k^{1+\varepsilon}$ tends to infinity, contradicting Conjecture~24.
\end{proof}
\section{Formal verification and scope}
The companion \href{https://github.com/Atm4x/my-ai-math-results/tree/main/pratt-conjecture-24/lean}{Lean development} (Lean 4.34.0, standard library only) checks all three ETF conditions, the digit filter, the carry argument, the induced relation and cyclic embedding. For the formal proof, it is convenient to avoid Stirling's formula altogether. Grouping edge words by their digit sums and weight sums, a finite pigeonhole argument gives the all-length estimate
\[
8^n\le (6n+1)^6\val(\Z_{3(7^n-1)+1}).
\]
Its principal theorem, \texttt{Pratt.not\_ValPrattConjecture24}, negates the necessary natural-power formulation: for every natural $K,N_0$ there is $N\ge N_0$, $N>0$, with $\val(\Z_N)^{100}>K N^{101}$. If the real-exponent conjecture held, choosing $\varepsilon=1/100$ and raising its bound to the 100th power would contradict this statement. This last elementary implication is given in prose, not formalized over the reals. The recorded kernel audit has only the standard axioms \texttt{propext}, \texttt{Classical.choice}, and \texttt{Quot.sound}.

This preprint does not settle the stronger SDPP conjecture or the matrix multiplication exponent. The sharper multinomial estimate is the informal argument above; the formal proof uses the alternative bound. Independent external review and scientific priority verification remain pending.
\section*{AI assistance and contributions}
The counterexample presented in this paper, its mathematical proof, and the accompanying Lean formalization were generated by ChatGPT (OpenAI). Artem Dikarev directed the research, selected the problem, guided the successive investigations, and initiated the preparation of this manuscript. ChatGPT also generated the verification scripts and assisted in writing and revising the exposition. The formal proof was checked with Lean 4.34.0 at trust level zero. Independent mathematical review and verification of scientific priority remain pending.
\begin{thebibliography}{1}
\bibitem{pratt} K. Pratt, On Generalized Corners and Matrix Multiplication, \emph{15th Innovations in Theoretical Computer Science Conference (ITCS 2024)}, LIPIcs 287, article 89, 2024. \href{https://doi.org/10.4230/LIPIcs.ITCS.2024.89}{doi:10.4230/LIPIcs.ITCS.2024.89}.
\end{thebibliography}
\end{document}
